\documentclass[a4paper,12pt,twoside]{memoir}

%otros
\usepackage{listings}
\interfootnotelinepenalty=10000
\usepackage{makecell}
\usepackage{listings}
\usepackage{dirtree}

% Castellano
\usepackage[spanish,es-tabla]{babel}
\selectlanguage{spanish}
\usepackage[utf8]{inputenc}
\usepackage[T1]{fontenc}
\usepackage{lmodern} % Scalable font
\usepackage{microtype}
\usepackage{placeins}

\RequirePackage{booktabs}
\RequirePackage[table]{xcolor}
\RequirePackage{xtab}
\RequirePackage{multirow} 

% Links
\usepackage[colorlinks]{hyperref}
\hypersetup{
	allcolors = {blue}
}

% Ecuaciones
\usepackage{amsmath}
\usepackage{multirow}

% Rutas de fichero / paquete
\newcommand{\ruta}[1]{{\sffamily #1}}

% Párrafos
\nonzeroparskip

% Huérfanas y viudas
\widowpenalty100000
\clubpenalty100000

% Evitar solapes en el header
\nouppercaseheads

% Imagenes
\usepackage{graphicx}
\newcommand{\imagen}[4]{
	\begin{figure}[!h]
		\centering
		\includegraphics[width=0.9\textwidth]{#1}
		\caption[{#3}]{#2}\label{#4}
	\end{figure}
	\FloatBarrier
}

\newcommand{\imagenflotante}[2]{
	\begin{figure}%[!h]
		\centering
		\includegraphics[width=0.9\textwidth]{#1}
		\caption{#2}\label{fig:#1}
	\end{figure}
}



% El comando \figura nos permite insertar figuras comodamente, y utilizando
% siempre el mismo formato. Los parametros son:
% 1 -> Porcentaje del ancho de página que ocupará la figura (de 0 a 1)
% 2 --> Fichero de la imagen
% 3 --> Texto a pie de imagen
% 4 --> Etiqueta (label) para referencias
% 5 --> Opciones que queramos pasarle al \includegraphics
% 6 --> Opciones de posicionamiento a pasarle a \begin{figure}
\newcommand{\figuraConPosicion}[6]{%
  \setlength{\anchoFloat}{#1\textwidth}%
  \addtolength{\anchoFloat}{-4\fboxsep}%
  \setlength{\anchoFigura}{\anchoFloat}%
  \begin{figure}[#6]
    \begin{center}%
      \Ovalbox{%
        \begin{minipage}{\anchoFloat}%
          \begin{center}%
            \includegraphics[width=\anchoFigura,#5]{#2}%
            \caption{#3}%
            \label{#4}%
          \end{center}%
        \end{minipage}
      }%
    \end{center}%
  \end{figure}%
}

%
% Comando para incluir imágenes en formato apaisado (sin marco).
\newcommand{\figuraApaisadaSinMarco}[5]{%
  \begin{figure}%
    \begin{center}%
    \includegraphics[angle=90,height=#1\textheight,#5]{#2}%
    \caption{#3}%
    \label{#4}%
    \end{center}%
  \end{figure}%
}
% Para las tablas
\newcommand{\otoprule}{\midrule [\heavyrulewidth]}
%
% Nuevo comando para tablas pequeñas (menos de una página).
\newcommand{\tablaSmall}[5]{%
 \begin{table}[h!]
  \begin{center}
   \rowcolors {2}{gray!35}{}
   \begin{tabular}{#2}
    \toprule
    #4
    \otoprule
    #5
    \bottomrule
   \end{tabular}
   \caption{#1}
   \label{tabla:#3}
  \end{center}
 \end{table}
}

%
% Nuevo comando para tablas pequeñas (menos de una página).
\newcommand{\tablaSmallSinColores}[5]{%
 \begin{table}[h!]
  \begin{center}
   \begin{tabular}{#2}
    \toprule
    #4
    \otoprule
    #5
    \bottomrule
   \end{tabular}
   \caption{#1}
   \label{tabla:#3}
  \end{center}
 \end{table}
}

\newcommand{\tablaApaisadaSmall}[5]{%
\begin{landscape}
  \begin{table}
   \begin{center}
    \rowcolors {2}{gray!35}{}
    \begin{tabular}{#2}
     \toprule
     #4
     \otoprule
     #5
     \bottomrule
    \end{tabular}
    \caption{#1}
    \label{tabla:#3}
   \end{center}
  \end{table}
\end{landscape}
}

%
% Nuevo comando para tablas grandes con cabecera y filas alternas coloreadas en gris.
\newcommand{\tabla}[6]{%
  \begin{center}
    \tablefirsthead{
      \toprule
      #5
      \otoprule
    }
    \tablehead{
      \multicolumn{#3}{l}{\small\sl continúa desde la página anterior}\\
      \toprule
      #5
      \otoprule
    }
    \tabletail{
      \hline
      \multicolumn{#3}{r}{\small\sl continúa en la página siguiente}\\
    }
    \tablelasttail{
      \hline
    }
    \bottomcaption{#1}
    \rowcolors {2}{gray!35}{}
    \begin{xtabular}{#2}
      #6
      \bottomrule
    \end{xtabular}
    \label{tabla:#4}
  \end{center}
}

%
% Nuevo comando para tablas grandes con cabecera.
\newcommand{\tablaSinColores}[6]{%
  \begin{center}
    \tablefirsthead{
      \toprule
      #5
      \otoprule
    }
    \tablehead{
      \multicolumn{#3}{l}{\small\sl continúa desde la página anterior}\\
      \toprule
      #5
      \otoprule
    }
    \tabletail{
      \hline
      \multicolumn{#3}{r}{\small\sl continúa en la página siguiente}\\
    }
    \tablelasttail{
      \hline
    }
    \bottomcaption{#1}
    \begin{xtabular}{#2}
      #6
      \bottomrule
    \end{xtabular}
    \label{tabla:#4}
  \end{center}
}

%
% Nuevo comando para tablas grandes sin cabecera.
\newcommand{\tablaSinCabecera}[5]{%
  \begin{center}
    \tablefirsthead{
      \toprule
    }
    \tablehead{
      \multicolumn{#3}{l}{\small\sl continúa desde la página anterior}\\
      \hline
    }
    \tabletail{
      \hline
      \multicolumn{#3}{r}{\small\sl continúa en la página siguiente}\\
    }
    \tablelasttail{
      \hline
    }
    \bottomcaption{#1}
  \begin{xtabular}{#2}
    #5
   \bottomrule
  \end{xtabular}
  \label{tabla:#4}
  \end{center}
}



\definecolor{cgoLight}{HTML}{EEEEEE}
\definecolor{cgoExtralight}{HTML}{FFFFFF}

%
% Nuevo comando para tablas grandes sin cabecera.
\newcommand{\tablaSinCabeceraConBandas}[5]{%
  \begin{center}
    \tablefirsthead{
      \toprule
    }
    \tablehead{
      \multicolumn{#3}{l}{\small\sl continúa desde la página anterior}\\
      \hline
    }
    \tabletail{
      \hline
      \multicolumn{#3}{r}{\small\sl continúa en la página siguiente}\\
    }
    \tablelasttail{
      \hline
    }
    \bottomcaption{#1}
    \rowcolors[]{1}{cgoExtralight}{cgoLight}

  \begin{xtabular}{#2}
    #5
   \bottomrule
  \end{xtabular}
  \label{tabla:#4}
  \end{center}
}


\graphicspath{ {./img/} }

% Capítulos
\chapterstyle{bianchi}
\newcommand{\capitulo}[2]{
	\setcounter{chapter}{#1}
	\setcounter{section}{0}
	\chapter*{#2}
	\addcontentsline{toc}{chapter}{#1. #2}
	\markboth{#2}{#2}
}

% Apéndices
\renewcommand{\appendixname}{Apéndice}
\renewcommand*\cftappendixname{\appendixname}

\newcommand{\apendice}[1]{
	%\renewcommand{\thechapter}{A}
	\chapter{#1}
}

\renewcommand*\cftappendixname{\appendixname\ }

% Formato de portada
\makeatletter
\usepackage{xcolor}
\newcommand{\tutor}[1]{\def\@tutor{#1}}
\newcommand{\course}[1]{\def\@course{#1}}
\definecolor{cpardoBox}{HTML}{E6E6FF}
\def\maketitle{
  \null
  \thispagestyle{empty}
  % Cabecera ----------------
\begin{center}%
	{\noindent\Huge Universidades de Burgos, León y Valladolid}\vspace{.5cm}%
	
	{\noindent\Large Máster universitario}\vspace{.5cm}%
	
	{\noindent\Huge \textbf{Inteligencia de Negocio y Big~Data en Entornos Seguros}}\vspace{.5cm}%
\end{center}%

\begin{center}%
	\includegraphics[height=3cm]{img/escudoUBU} \hspace{1cm}
	\includegraphics[height=3cm]{img/escudoUVA} \hspace{1cm}
	\includegraphics[height=3cm]{img/escudoULE} \vspace{1cm}%
\end{center}%

  \vfill
  % Título proyecto y escudo informática ----------------
  \colorbox{cpardoBox}{%
    \begin{minipage}{.9\textwidth}
      \vspace{.5cm}\Large
      \begin{center}
      \textbf{TFM del Máster Inteligencia de Negocio y Big Data en Entornos Seguros}\vspace{.6cm}\\
      \textbf{\LARGE\@title{}}
      \end{center}
      \vspace{.2cm}
    \end{minipage}

  }%
  \hfill
  \vfill
  % Datos de alumno, curso y tutores ------------------
  \begin{center}%
  {%
    \noindent\LARGE
    Presentado por \@author{}\\ 
    en Universidad de Burgos --- \@date{}\\
    Tutores: \@tutor{}\\
  }%
  \end{center}%
  \null
  \cleardoublepage
  }
\makeatother

\newcommand{\nombre}{Ignacio Ponsoda Llorens} %%% cambio de comando

% Datos de portada
\title{Clasificación de individuos a partir de imágenes oculares y redes neuronales pre-entrenadas.}
\author{\nombre}
\tutor{Dr. José Francisco Díez Pastor y Dr. Pedro Latorre Carmona}
\date{\today}

\begin{document}

\maketitle


\newpage\null\thispagestyle{empty}\newpage


%%%%%%%%%%%%%%%%%%%%%%%%%%%%%%%%%%%%%%%%%%%%%%%%%%%%%%%%%%%%%%%%%%%%%%%%%%%%%%%%%%%%%%%%
\thispagestyle{empty}


\noindent
\begin{center}%
	{\noindent\Huge Universidades de Burgos, León y Valladolid}\vspace{.5cm}%
	
\begin{center}%
	\includegraphics[height=3cm]{img/escudoUBU} \hspace{1cm}
	\includegraphics[height=3cm]{img/escudoUVA} \hspace{1cm}
	\includegraphics[height=3cm]{img/escudoULE} \vspace{1cm}%
\end{center}%

	{\noindent\Large \textbf{Máster universitario en Inteligencia de Negocio y Big~Data en Entornos Seguros}}\vspace{.5cm}%
\end{center}%



\noindent D. José Francisco Díez Pastor, profesor del departamento de Ingeniería Informática, área de Lenguajes y Sistemas Informáticos, y
\noindent D. Pedro Latorre Carmona, profesor del departamento de Ingeniería Informática, área de Lenguajes y Sistemas Informáticos.

\noindent Exponen:

\noindent Que el alumno  \nombre, con DNI 21698927Z, ha realizado el Trabajo final de Máster en Inteligencia de Negocio y Big Data en Entornos Seguros 
          titulado "Clasificación de indivíduos a partir de imágenes oculares y redes neuronales pre-entrenadas". 

\noindent Y que dicho trabajo ha sido realizado por el alumno bajo la dirección de los que suscriben, en virtud de lo cual se autoriza su presentación y defensa.

\begin{center} %\large
En Burgos, {\large \today}
\end{center}

\vfill\vfill\vfill

% Author and supervisor
\begin{minipage}{0.45\textwidth}
\begin{flushleft} %\large
Vº. Bº. del Tutor:\\[2cm]
D. José Francisco Díez Pastor
\end{flushleft}
\end{minipage}
\hfill
\begin{minipage}{0.45\textwidth}
\begin{flushleft} %\large
Vº. Bº. del tutor:\\[2cm]
D. Pedro Latorre Carmona
\end{flushleft}
\end{minipage}
\hfill

\vfill

% para casos con solo un tutor comentar lo anterior
% y descomentar lo siguiente
%Vº. Bº. del Tutor:\\[2cm]
%D. nombre tutor


\newpage\null\thispagestyle{empty}\newpage




\frontmatter

% Abstract en castellano
\renewcommand*\abstractname{Resumen}
\begin{abstract}
La utilización de la biometría para mejorar la seguridad, principalmente en lo referente al acceso de dispositivos electrónicos, es un recurso ampliamente empleado en la actualidad. El iris uno de los elementos biométricos que mayores dificultades presentan para su suplantación, y es por ello que su utilización en este campo ha atraído la atención de la comunidad científica estas últimas dos décadas. 

En este proyecto se han adaptado redes neuronales, inicialmente entrenadas para clasificar diversos objetos, para que sean capaces de identificar a un individuo utilizando su imagen ocular. 

Para ello, se han utilizado dos enfoques. En un primer enfoque, las redes neuronales se ha adaptado utilizando imágenes oculares completas, mientras que para el segundo enfoque, se ha hecho lo propio, pero aislado la zona del iris, que es apriori la zona de la imagen ocular que mejor permite la identificación de individuos.

Además, se han utilizado técnicas de ampliación del \textit{dataset} original, a fin de contar con un mayor número de muestras de cada indivíduo y también, mejorar la robustez de las redes neuronales adaptadas.

Los resultados muestran que las mejores tasas de clasificación se han dado en el enfoque donde se utilizaba la imagen ocular completa, sin que las técnicas de ampliación del \textit{dataset} hayan permitido mejorar la tasa de clasificación. 

Como futuras líneas de trabajo, se establecen la utilización de redes neuronales pre-entrenadas distintas así como testear el modelo con imágenes realizadas fuera del entorno académico.


\end{abstract}

\renewcommand*\abstractname{Descriptores}
\begin{abstract}
biometría, iris, redes neuronales
\end{abstract}

\clearpage

% Abstract en inglés 
\renewcommand*\abstractname{Abstract}
\begin{abstract}
  The use of biometrics to improve the security of electronic devices is a widely used resource nowadays. The iris is one of the biometric elements available to the human being.
  This work analyzes the recognition capacity of individuals of adapted neural networks. To do this, using fine-tuning, four variations of the dataset have been adapted to the VGG16 network.
  Two of these variations correspond, on the one hand, to the dataset without preprocessing, and on the other, to images where the iris has been segmented and normalized.
  The two remaining variations have been determined by the use of data augmentation on the two previously described.

  The results show that the neural network has worked more efficiently in the images where the iris area had not been isolated.
  Likewise, the network has had a better classification capacity when the data augmentation has not been applied, something that can
  This may be because the images have not been tested in real scenarios. 

  Based on these results, it is possible to extract the capacity of the neural networks themselves to find features of the human eye beyond the iris, and thus improve their classification capacity.
\end{abstract}

\renewcommand*\abstractname{Keywords}
\begin{abstract}
biometrics, iris, neural networks, fine-tuning, data augmentation
\end{abstract}

\clearpage

% Indices
\tableofcontents

\clearpage

\listoffigures

\clearpage

\listoftables
\clearpage

\mainmatter

\addcontentsline{toc}{part}{Memoria}
\part*{Memoria}

\capitulo{1}{Introducción}

El iris se utiliza como elemento de reconocimiento biométrico, tanto por su inmutabilidad a lo largo del tiempo como por resultar un valor único y 
personal, que supone que dos personas no tendrían el mismo iris (05 Iris Recognition Developmen Techniques: A Comprehensive Review).

En este estudio, se compararán técnicas de redes neuronales completas como una combinación de redes neuronales para la segmentación del ojo, la extracción
de características y finalmente, con técnicas de Machine Learning.

Para ello, vamos a utilizar una red neuronal preentrenada de ImageNet a la que le aplicaremos fine-tuning y compararemos los resultados con la combinación de
una red neuronal preentrenada para la segmentación, un preprocesamiento para extraer y normalizar el iris, la red preentrenada ImageNet para extraer las características principales
(sin realizar fine-tuning) y machine learning para el cluster (identificación del individuo).

\section{Outline}

El resto del docuemnto se estructura de la siguiente manera. El capítulo 2 contiene los objetivos del proyecto. El capítulo 3 contiene los conceptos teóricos necesarios
para entender el proyecto. El capítulo 4 muestra las técnicas y herramientas utilizadas en el desarrollo de este trabajo. En el capítulo 5 se muestran los Aspectos
más relevantes que se han desarrollado. En el capítulo 6 los trabajos relacionados y en el capítulo 7 las conclusiones y las líneas de trabajo futuras.

\capitulo{2}{Objetivos del proyecto}

Este apartado explica de forma precisa y concisa cuales son los objetivos que se persiguen con la realización del proyecto. Se puede distinguir entre los objetivos marcados por los requisitos del software a construir y los objetivos de carácter técnico que plantea a la hora de llevar a la práctica el proyecto.

\capitulo{3}{Conceptos teóricos}

05 techniques review entre todos los elementos biometricos, como huellas, cara, iris, voz y , el iris recognition system (IRS) es el metodo con mayor eficiencia a la hora de determinar la identiidad
de las personas, ya que el iris es el mismo a lo largo de la vida de una persona y es úncio, incluso entre gemelos. Esto es incluso utilizado en los procesos 
post-mortem para poder determinar la pertenencia del cuerpo 06 Post-mortem iris recognition. El iris humanos es un organo del ojo, que tiene como funciona controlar el tamaño
de la pupila en función de la cantidad de luz que llega hasta este 06 Post-mortem iris recognition.

\section{Secciones}

Las secciones se incluyen con el comando section.

\subsection{Subsecciones}

Además de secciones tenemos subsecciones.

\subsubsection{Subsubsecciones}

Y subsecciones. 


\section{Referencias}

Las referencias se incluyen en el texto usando cite \cite{wiki:latex}. Para citar webs, artículos o libros \cite{koza92}.


\section{Imágenes}

Se pueden incluir imágenes con los comandos standard de \LaTeX, pero esta plantilla dispone de comandos propios como por ejemplo el siguiente:


\section{Listas de items}

Existen tres posibilidades:

\begin{itemize}
	\item primer item.
	\item segundo item.
\end{itemize}

\begin{enumerate}
	\item primer item.
	\item segundo item.
\end{enumerate}

\begin{description}
	\item[Primer item] más información sobre el primer item.
	\item[Segundo item] más información sobre el segundo item.
\end{description}
	
\begin{itemize}
\item 
\end{itemize}

\section{Tablas}

Igualmente se pueden usar los comandos específicos de \LaTeX o bien usar alguno de los comandos de la plantilla.

\tablaSmall{Herramientas y tecnologías utilizadas en cada parte del proyecto}{l c c c c}{herramientasportipodeuso}
{ \multicolumn{1}{l}{Herramientas} & App AngularJS & API REST & BD & Memoria \\}{ 
HTML5 & X & & &\\
CSS3 & X & & &\\
BOOTSTRAP & X & & &\\
JavaScript & X & & &\\
AngularJS & X & & &\\
Bower & X & & &\\
PHP & & X & &\\
Karma + Jasmine & X & & &\\
Slim framework & & X & &\\
Idiorm & & X & &\\
Composer & & X & &\\
JSON & X & X & &\\
PhpStorm & X & X & &\\
MySQL & & & X &\\
PhpMyAdmin & & & X &\\
Git + BitBucket & X & X & X & X\\
Mik\TeX{} & & & & X\\
\TeX{}Maker & & & & X\\
Astah & & & & X\\
Balsamiq Mockups & X & & &\\
VersionOne & X & X & X & X\\
} 

\capitulo{4}{Técnicas y herramientas}

Explicación de la sección

\section{Dataset}

Se ha utilizado el dataset de CASIA para la imágenes de ojos.

\section{Redes neuronales preentrenadas}

Así mismo, se ha utilizado la red neuronal preentrenada Basada en U-Net y accesible desde https://github.com/jus390/U-net-Iris-segmentation, la cual ya había sido
entrenada para la segmentación del iris.

Finalmente, para el último proceso del proyecto, se utiliza imagenet como red neuronal preentrenada central.

\section{Github}
Para el control de versiones y el seguimiento de fases del proyecto

\section{Python}

Además se ha utilizado python para todo el proyecto, esto quiere decir preprocesado, creación y utilización de redes neuronales y clasificación contécnicas de 
machine learning.

Entre las principales librerías utilizadas se encuentran:

\begin{itemize}
    \item os, para el acceso a los directorios
    \item numpy, para trabajar con las imágenes a nivel de arrays
    \item scikit-image, para la transformación de las imágenes y el uso de dataset
    \item tensorflow, para la modificación de las redes neuronales
    \item keras, para el manejo de las redes neuronales
    \item matplot, para las gráficas
\end{itemize}

\subsection{Visual Studio}

Tanto para la redacción de la memoria con latex como para el estudio.

\section{Scrum}
Scrum es un marco que ayuda a la organización de los equipos entorno a un proyecto, en base a Sprints de una determinada duración y permite una retroalimentación
contínua del proyecto, de forma que se asegura que el equipo trabaje en consonancia.

En el caso de este proyecto, se ha adecuado la metodología de Scrum para cuadrar reuniones semanales o bisemanales con los tutores del proyecto, estableciendo
los sprints y los objetivos de cada sprint utilizando la plataforma Github.

\section{Recursos}

En cuanto a las herramientas que se han contado para realizar los entrenamientos de las redes neuronales, se ha contado con un ordenador portatil con un procesador
i7, 12 gb de ram y una gráfica dedicada NVIDIA de 1 gb.





\capitulo{5}{Aspectos relevantes del desarrollo del proyecto}

El proyecto puede dividirse en una fase previa y tres propuestas diferentes para la identificación biomédica de personas a través de una imagen de su ojo.

La primera fase trata de la optimización del código previo y la creación de un pipeline con una configuración que permitiese un mejor control de las funciones. 
También del  data augnmentation, en el cual se aplica tanto ruido gausseano (de 2.5, 5 y 7.5) como transformaciones afines (), siendo realizadas de 
forma aleatoria, con lo cual los supuestos pueden ser a) imagen sin data augnmentation, b)imagen con ruido gausseano, c)imagen con transformaciones afines y 
d)imagen con ruido gausseano y transformaciones afines. Este dataset será el base para todo el proyecto.

La segunda fase se trata de la elección de la mejor forma de clasificar las imagenes de entrada:

1.La propuesta en el TFG anterior, en la cual se realiza un preprocesamiento de las imágenes del dataset original. Este preprocesamiento consiste primer lugar 
en la segmentación de las imágenes del iris con una red neuronal preentrenada precisamente para realizar esta acción. En segundo lugar, se realiza una extracción
del iris a través de una binarización de las partes del ojo y una extracción del iris, a la que se le aplica una normalización para que quede proyectado.
La siguiente fase de esta primera propuesta es la extracción de características (quitandole las dos últimas capas a una red neuronal) de la imagen normalizada
 para posteriormente utilizar una red neuronal preentrenada con imagenet (de hecho 3 redes de la cual se elige la mejor). Posteriormente, estas características
 extraidas se pasan a modelos de ML, que son los que realizarán la clasificación. 

 Todo el proyecto queda establecido en una pipeline

\capitulo{6}{Trabajos relacionados}

El principal apartado anterior lo podemos encontrar en el TFG de extracción del iris, en el cual se basa este trarbajo, puesto se realiza a grandes rasgos
todo lo relativo a la primera de las opciones del trabajo.

Sobre temas de extracción del iris encontramos 02 iris wavelet neural, donde se hace un preprocesamiento con extración del iris utilizando Hough Transform y la 
normalización con Daugmands rubber. Luego, tras eliminar el ruido, la extración se realiza con transformaciones de wavelet. Finalmente, se crea una red neutornal
utilizando el mean-squared error para calcular los pesos en la red.

En 03 deep iris encontramos el desarrollo de técnicas de deep learning para el reconocimiento del iris basado en convolutional neural network residual. utilizando una red preentrenada 
de ResNet50 y fine-tuning, entrenado con una cross-entropy loss function (but they are not using data augnmentation, they are using another dataset IIT Delhi dataset and they are 
not doing the p`reprocessing step).

In 13 ImageNet Deep CNN they use the ImageNet dataset with data augnmentation, dropout to train a nerual network to detect images of the feed (maybe we can remove it as
it is not related with iris).

En 14 experimental deep convolutional iris recognition, utilizaron el dataset CASIA - 10000 y la arquitectura VGG-Net, lo cual realiza un PCA para extraer los elementos
más característicos de las imágenes . Después utilizan algoritmos de clasificación para clasificar las imágenes, como el SVM (esto es similar al TFG) y consiguen unos percentajes
de reconocimiento muy altos.

05 techniques review habla de siete pasos en los que se divide un sistema de reconocimiento del iris, 1)adquisición, 2)preprocesamiento, 3)segmentación
 4)normalización, 5)extracción de características, 6) selección de features únicos y característicos, 7)clasificación. Este paper también describe una falta de
 trabajos entorno a datasets de baja cualidad (revisar para sacar más papers). El paper también realza que los sistemas de reconocimiento del iris (IRS) se vuelven 
 poco efectivos cuando las imagenes tienen rotaciones or reflejos, algo que intentamos de mejorar en nuestro proceso, añadiendo ruido con el data augnmentation.

 Este mismo paper también comenta los distintos dataset utilizados para estos estudios de reconocimiento de iris, el tipo de ruido utilizado así como su método,
 los tipos de segmentación tradicional y actual utilizados (habitualmente con redes neuronales), técnicas de normalización y extracciñon de caracterñisticas, así 
 como los tipos de accuracy de los métodos de iris recognition.
.... (utilizando Machine learning y deep learning)

Y para el tema de clasificación de personas con el ojo, encontramos ...

Por otro lado, en temas de fine tunning con redes neuronales, tenemos...
\capitulo{7}{Conclusiones y Líneas de trabajo futuras}

Todo proyecto debe incluir las conclusiones que se derivan de su desarrollo. Éstas pueden ser de diferente índole, dependiendo de la tipología del proyecto, pero normalmente van a estar presentes un conjunto de conclusiones relacionadas con los resultados del proyecto y un conjunto de conclusiones técnicas. 
Además, resulta muy útil realizar un informe crítico indicando cómo se puede mejorar el proyecto, o cómo se puede continuar trabajando en la línea del proyecto realizado. 



%\renewcommand\chaptername{Anexo}
%\renewcommand\thechapter{\Roman{chapter}}
%\setcounter{chapter}{0}

% Añadir entrada en el índice: Anexos
\appendix
\addcontentsline{toc}{part}{Apéndices}
\part*{Apéndices}

\apendice{Plan de Proyecto Software}

\section{Introducción}

\section{Planificación temporal}

\section{Estudio de viabilidad}

\subsection{Viabilidad económica}

\subsection{Viabilidad legal}



\apendice{Especificación de Requisitos}

\section{Introducción}

\section{Objetivos generales}

\section{Catalogo de requisitos}

\section{Especificación de requisitos}



\apendice{Especificación de diseño}

\section{Introducción}

\section{Diseño de datos}

\section{Diseño procedimental}

\section{Diseño arquitectónico}



\apendice{Documentación técnica de programación}

\section{Introducción}

\section{Estructura de directorios}

\section{Manual del programador}

\section{Compilación, instalación y ejecución del proyecto}

\section{Pruebas del sistema}

%\apendice{Documentación de usuario}

\section{Introducción}

\section{Requisitos de usuarios}

\section{Instalación}

\section{Manual del usuario}





\bibliographystyle{plain}
\bibliography{bibliografia}

\end{document}
